Averaging over all combinations of the other two parameters, propagation pressure produced the largest broad deterioration in containment outcomes.
Increasing \(p\) from 0.001 to 0.020 changed mean systemic contagion from \(0.01\%\) to \(8.94\%\), a difference of \(8.92\) percentage points, while cross-cell leakage changed from \(2.14\%\) to \(25.64\%\).
Increasing mediation delay from \(\tau=0\) to \(\tau=6\) changed mean systemic contagion from \(0.91\%\) to \(6.50\%\) and reduced mean strategic velocity from \(99.33\%\) to \(95.60\%\).
Increasing the false-negative rate from \(0\%\) to \(30\%\) changed mean systemic contagion from \(2.03\%\) to \(4.21\%\) and mean cross-cell leakage from \(9.68\%\) to \(16.27\%\).
The value-related outcomes follow the same structural pattern but capture different mechanisms.
Across the low-to-high \(p\) contrast, propagated value-missing exposure changed by \(0.359\) secondary-cell equivalents and value slippage by \(1.452\) compromised cell-steps; the corresponding delay contrasts were \(0.241\) and \(6.629\), and the detection-failure contrasts were \(0.096\) and \(3.145\).
These contrasts are descriptive marginal comparisons over the specified experimental grid rather than empirical effect estimates.
